% GENERATED FILE. Edit canonical lessons, then run npm run generate:books.
% canonical-source-hash: fnv1a64:17b75aa01c70a127
% canonical-lessons: KA-C478-haddu, KA-C478-kokkare, KA-C478-mincuhulu, KA-C478-edi, KA-C478-timingila

\chapter{From Eagles to Whales}
\label{ch:ka-from-eagles-to-whales}
\hlchaptermodality{478}

\begin{chapteropening}
\textbf{By the end of this chapter:} \emph{I can say an eagle, a stork, a firefly, a crab and a whale.}

Complete the last lesson of chapter 478: I can say an eagle, a stork, a firefly, a crab and a whale.
\end{chapteropening}

\section[\texorpdfstring{\kn{ಹದ್ದು} (haddu)}{ಹದ್ದು (haddu)}]{\kn{ಹದ್ದು} (haddu) — an eagle; a kite}

\label{lesson:KA-C478-haddu}



\begin{quote}
Before the new one: say the Kannada for to graze, then the Kannada for a termite.
\end{quote}

\subsection*{You'll want to know: \kn{ಹದ್ದು}}
\textbf{\kn{ಹದ್ದು}} — \emph{haddu} — \textquotedblleft{}an eagle; a kite\textquotedblright{}.

\begin{quote}
\kn{ಆಕಾಶದಲ್ಲಿ} \kn{ಒಂದು} \kn{ಹದ್ದು} \kn{ಸುತ್ತುತ್ತಿತ್ತು}.
\end{quote}

\begin{quote}
\emph{ākāśadalli ondu haddu suttuttittu.}
\end{quote}

\begin{quote}
An eagle was circling in the sky.
\end{quote}

A large bird of prey that circles high up.

\begin{cousinweb}
\textbf{\kn{ಹದ್ದು}} (\emph{haddu}) is native Dravidian: DEDR 3977 lists Kannada \emph{pardu}, \emph{paddu}, \textquotedblleft{}a kite, a hawk, a vulture\textquotedblright{}, the old \emph{p-} turned into \emph{h-}.
\end{cousinweb}

\begin{grammarlens}[title={What you've built}]
One more word for getting around.
\end{grammarlens}

\subsection*{Your turn}
\begin{itemize}
\raggedright
  \item \emph{Say it:} \emph{haddu}
  \item \emph{Say it:} \emph{haddu}, once more
  \item \emph{Say it:} \emph{ākāśadalli ondu haddu suttuttittu.} — An eagle was circling in the sky.
  \item \emph{Recall:} say the Kannada for to graze, then the Kannada for a termite, then say \emph{haddu} again
\end{itemize}

\begin{tcolorbox}[breakable,colback=teal!4,colframe=teal!35!black,title={Before you move on}]
What does \kn{ಹದ್ದು} mean? (\textquotedblleft{}An eagle; a kite\textquotedblright{}.) Say it once more.
\end{tcolorbox}
\section[\texorpdfstring{\kn{ಕೊಕ್ಕರೆ} (kokkare)}{ಕೊಕ್ಕರೆ (kokkare)}]{\kn{ಕೊಕ್ಕರೆ} (kokkare) — a stork; a crane}

\label{lesson:KA-C478-kokkare}



\begin{quote}
Before the new one: say the Kannada for a termite, then the Kannada for an eagle.
\end{quote}

\subsection*{You'll want to know: \kn{ಕೊಕ್ಕರೆ}}
\textbf{\kn{ಕೊಕ್ಕರೆ}} — \emph{kokkare} — \textquotedblleft{}a stork; a crane\textquotedblright{}.

\begin{quote}
\kn{ಕೆರೆಯ} \kn{ದಡದಲ್ಲಿ} \kn{ಬಿಳಿ} \kn{ಕೊಕ್ಕರೆಗಳು} \kn{ನಿಂತಿದ್ದವು}.
\end{quote}

\begin{quote}
\emph{kereya daḍadalli biḷi kokkaregaḷu nintiddavu.}
\end{quote}

\begin{quote}
White storks were standing on the bank of the lake.
\end{quote}

Plural: \kn{ಕೊಕ್ಕರೆಗಳು}, \emph{kokkaregaḷu}.

A long-legged bird of fields and lakes.

\begin{cousinweb}
\textbf{\kn{ಕೊಕ್ಕರೆ}} (\emph{kokkare}) is native Dravidian: DEDR 2125 lists Kannada \emph{kokku}, \emph{kokkare}, \textquotedblleft{}a crane\textquotedblright{}.
\end{cousinweb}

\begin{grammarlens}[title={What you've built}]
One more word for getting around.
\end{grammarlens}

\subsection*{Your turn}
\begin{itemize}
\raggedright
  \item \emph{Say it:} \emph{kokkare}
  \item \emph{Say it:} \emph{kokkare}, once more
  \item \emph{Say it:} \emph{kereya daḍadalli biḷi kokkaregaḷu nintiddavu.} — White storks were standing on the bank of the lake.
  \item \emph{Recall:} say the Kannada for a termite, then the Kannada for an eagle, then say \emph{kokkare} again
\end{itemize}

\begin{tcolorbox}[breakable,colback=teal!4,colframe=teal!35!black,title={Before you move on}]
What does \kn{ಕೊಕ್ಕರೆ} mean? (\textquotedblleft{}A stork; a crane\textquotedblright{}.) Say it once more.
\end{tcolorbox}
\section[\texorpdfstring{\kn{ಮಿಂಚುಹುಳು} (miñcuhuḷu)}{ಮಿಂಚುಹುಳು (miñcuhuḷu)}]{\kn{ಮಿಂಚುಹುಳು} (miñcuhuḷu) — a firefly}

\label{lesson:KA-C478-mincuhulu}



\begin{quote}
Before the new one: say the Kannada for an eagle, then the Kannada for a stork.
\end{quote}

\subsection*{You'll want to know: \kn{ಮಿಂಚುಹುಳು}}
\textbf{\kn{ಮಿಂಚುಹುಳು}} — \emph{miñcuhuḷu} — \textquotedblleft{}a firefly\textquotedblright{}.

\begin{quote}
\kn{ರಾತ್ರಿ} \kn{ಹೊಲದಲ್ಲಿ} \kn{ಮಿಂಚುಹುಳುಗಳು} \kn{ಹೊಳೆಯುತ್ತಿದ್ದವು}.
\end{quote}

\begin{quote}
\emph{rātri holadalli miñcuhuḷugaḷu hoḷeyuttiddavu.}
\end{quote}

\begin{quote}
At night fireflies were glowing in the field.
\end{quote}

\kn{ಮಿಂಚು}, lightning, plus \kn{ಹುಳು}, an insect or worm.

\begin{cousinweb}
\textbf{\kn{ಮಿಂಚುಹುಳು}} (\emph{miñcuhuḷu}) is a native Kannada compound of \kn{ಮಿಂಚು}, \textquotedblleft{}lightning, a flash\textquotedblright{} (DEDR 4876), and \kn{ಹುಳು}, \textquotedblleft{}a worm, an insect\textquotedblright{}.
\end{cousinweb}

\begin{grammarlens}[title={What you've built}]
One more word for getting around.
\end{grammarlens}

\subsection*{Your turn}
\begin{itemize}
\raggedright
  \item \emph{Say it:} \emph{miñcuhuḷu}
  \item \emph{Say it:} \emph{miñcuhuḷu}, once more
  \item \emph{Say it:} \emph{rātri holadalli miñcuhuḷugaḷu hoḷeyuttiddavu.} — At night fireflies were glowing in the field.
  \item \emph{Recall:} say the Kannada for an eagle, then the Kannada for a stork, then say \emph{miñcuhuḷu} again
\end{itemize}

\begin{tcolorbox}[breakable,colback=teal!4,colframe=teal!35!black,title={Before you move on}]
What does \kn{ಮಿಂಚುಹುಳು} mean? (\textquotedblleft{}A firefly\textquotedblright{}.) Say it once more.
\end{tcolorbox}
\section[\texorpdfstring{\kn{ಏಡಿ} (ēḍi)}{ಏಡಿ (ēḍi)}]{\kn{ಏಡಿ} (ēḍi) — a crab}

\label{lesson:KA-C478-edi}



\begin{quote}
Before the new one: say the Kannada for a stork, then the Kannada for a firefly.
\end{quote}

\subsection*{You'll want to know: \kn{ಏಡಿ}}
\textbf{\kn{ಏಡಿ}} — \emph{ēḍi} — \textquotedblleft{}a crab\textquotedblright{}.

\begin{quote}
\kn{ಕಡಲತೀರದ} \kn{ಮರಳಿನಲ್ಲಿ} \kn{ಏಡಿಗಳು} \kn{ಓಡುತ್ತಿದ್ದವು}.
\end{quote}

\begin{quote}
\emph{kaḍalatīrada maraḷinalli ēḍigaḷu ōḍuttiddavu.}
\end{quote}

\begin{quote}
Crabs were running about on the sand of the beach.
\end{quote}

Plural: \kn{ಏಡಿಗಳು}, \emph{ēḍigaḷu}.

The plural takes -\kn{ಗಳು}: \kn{ಏಡಿಗಳು}.

\begin{cousinweb}
\textbf{\kn{ಏಡಿ}} (\emph{ēḍi}) is native Dravidian: DEDR 2901 lists Kannada \emph{ēḍi}, beside \emph{esaḍi} and \emph{naḷḷi}, with Tamil \emph{ñaṇṭu}, \textquotedblleft{}a crab\textquotedblright{}.
\end{cousinweb}

\begin{grammarlens}[title={What you've built}]
One more word for getting around.
\end{grammarlens}

\subsection*{Your turn}
\begin{itemize}
\raggedright
  \item \emph{Say it:} \emph{ēḍi}
  \item \emph{Say it:} \emph{ēḍi}, once more
  \item \emph{Say it:} \emph{kaḍalatīrada maraḷinalli ēḍigaḷu ōḍuttiddavu.} — Crabs were running about on the sand of the beach.
  \item \emph{Recall:} say the Kannada for a stork, then the Kannada for a firefly, then say \emph{ēḍi} again
\end{itemize}

\begin{tcolorbox}[breakable,colback=teal!4,colframe=teal!35!black,title={Before you move on}]
What does \kn{ಏಡಿ} mean? (\textquotedblleft{}A crab\textquotedblright{}.) Say it once more.
\end{tcolorbox}
\section[\texorpdfstring{\kn{ತಿಮಿಂಗಿಲ} (timiṅgila)}{ತಿಮಿಂಗಿಲ (timiṅgila)}]{\kn{ತಿಮಿಂಗಿಲ} (timiṅgila) — a whale}

\label{lesson:KA-C478-timingila}



\begin{quote}
Before the new one: say the Kannada for a firefly, then the Kannada for a crab.
\end{quote}

\subsection*{You'll want to know: \kn{ತಿಮಿಂಗಿಲ}}
\textbf{\kn{ತಿಮಿಂಗಿಲ}} — \emph{timiṅgila} — \textquotedblleft{}a whale\textquotedblright{}.

\begin{quote}
\kn{ಸಮುದ್ರದಲ್ಲಿ} \kn{ನಾವು} \kn{ಒಂದು} \kn{ತಿಮಿಂಗಿಲವನ್ನು} \kn{ನೋಡಿದೆವು}.
\end{quote}

\begin{quote}
\emph{samudradalli nāvu ondu timiṅgilavannu nōḍidevu.}
\end{quote}

\begin{quote}
We saw a whale in the sea.
\end{quote}

In Sanskrit stories a huge fish that swallows other fish.

\begin{cousinweb}
\textbf{\kn{ತಿಮಿಂಗಿಲ}} (\emph{timiṅgila}) is borrowed from Sanskrit \emph{timiṃgila}, \textquotedblleft{}a large fabulous fish\textquotedblright{} (Monier-Williams), the swallower of the \emph{timi} fish.
\end{cousinweb}

\begin{grammarlens}[title={What you've built}]
One more word for getting around.
\end{grammarlens}

\subsection*{Your turn}
\begin{itemize}
\raggedright
  \item \emph{Say it:} \emph{timiṅgila}
  \item \emph{Say it:} \emph{timiṅgila}, once more
  \item \emph{Say it:} \emph{samudradalli nāvu ondu timiṅgilavannu nōḍidevu.} — We saw a whale in the sea.
  \item \emph{Recall:} say the Kannada for a firefly, then the Kannada for a crab, then say \emph{timiṅgila} again
\end{itemize}

\begin{tcolorbox}[breakable,colback=teal!4,colframe=teal!35!black,title={Before you move on}]
What does \kn{ತಿಮಿಂಗಿಲ} mean? (\textquotedblleft{}A whale\textquotedblright{}.) Say it once more.
\end{tcolorbox}
